\documentclass{article}
\usepackage{amsmath,amssymb,amsthm}
\usepackage{algorithm}
\usepackage{algorithmic}
\usepackage{bm}
\usepackage{hyperref}
\usepackage{enumitem}
\usepackage{graphicx}
\usepackage{xcolor}
\newtheorem{theorem}{Theorem}
\newtheorem{lemma}{Lemma}
\newtheorem{assumption}{Assumption}
\newcommand{\E}{\mathbb{E}}
\newcommand{\R}{\mathbb{R}}

\begin{document}

\section*{SignSGD (Sign of Stochastic Gradient Descent)}

\section*{Mathematical Formulation}
Let $f:\R^{d}\to\R$ be a (possibly non‑convex) loss function.  
At iteration $k$ a stochastic gradient $g_k\in\R^{d}$ is sampled such that
\[
g_k = \nabla f(x_k;\xi_k),
\]
where $\xi_k$ denotes the random data mini‑batch.  
Define the elementwise sign operator
\[
\operatorname{sign}(v)_i = 
\begin{cases}
+1 & v_i>0,\\
-1 & v_i<0,\\
0  & v_i=0 .
\end{cases}
\]
The update rule of **SignSGD** with learning rate $\alpha>0$ is
\begin{equation}
\boxed{\,x_{k+1}=x_k-\alpha\,\operatorname{sign}(g_k)\,}. \tag{1}
\end{equation}
Only the direction (sign) of each coordinate is transmitted; the magnitude is discarded.

\section*{Pseudocode}
\begin{algorithm}[H]
\caption{SignSGD}
\begin{algorithmic}[1]
\REQUIRE Initial point $x_0\in\R^{d}$, step size $\alpha>0$, number of iterations $T$.
\FOR{$k=0,\dots,T-1$}
    \STATE Sample a mini‑batch $\xi_k$ and compute stochastic gradient $g_k=\nabla f(x_k;\xi_k)$.
    \STATE $s_k \gets \operatorname{sign}(g_k)$ \COMMENT{elementwise sign}
    \STATE $x_{k+1} \gets x_k - \alpha\, s_k$.
\ENDFOR
\RETURN $x_T$ (or the averaged iterate $\bar x_T=\frac1T\sum_{k=0}^{T-1}x_k$).
\end{algorithmic}
\end{algorithm}

\section*{Dry Run Example}
Consider the one‑dimensional quadratic
\[
f(w)=(w-3)^2,
\]
with deterministic gradient $\nabla f(w)=2(w-3)$.  
Set $w_0=0$, learning rate $\alpha=0.1$, and run three iterations.

\begin{center}
\begin{tabular}{c|c|c|c}
Iteration $k$ & $w_k$ & $g_k=\nabla f(w_k)$ & $w_{k+1}=w_k-\alpha\,\operatorname{sign}(g_k)$\\\hline
0 & $0$ & $2(0-3)=-6$ & $0-0.1\cdot(-1)=0.1$\\
1 & $0.1$ & $2(0.1-3)=-5.8$ & $0.1-0.1\cdot(-1)=0.2$\\
2 & $0.2$ & $2(0.2-3)=-5.6$ & $0.2-0.1\cdot(-1)=0.3$
\end{tabular}
\end{center}
The objective values are
\[
f(0)=9,\qquad f(0.1)=8.41,\qquad f(0.2)=7.84,\qquad f(0.3)=7.29,
\]
showing monotone descent despite using only the sign information.

\newpage
\section*{Assumptions}
We work under the following standard conditions for non‑convex stochastic optimization.

\begin{assumption}[Smoothness]\label{ass:smooth}
$f$ is $L$‑smooth, i.e.
\[
\|\nabla f(x)-\nabla f(y)\|_2\le L\|x-y\|_2,\qquad\forall x,y\in\R^{d}.
\]
\end{assumption}

\begin{assumption}[Unbiased stochastic gradient]\label{ass:unbiased}
\[
\E_{\xi_k}[\,g_k\mid x_k\,]=\nabla f(x_k),\qquad\forall k.
\]
\end{assumption}

\begin{assumption}[Coordinate‑wise median bound]\label{ass:median}
For each coordinate $i\in\{1,\dots,d\}$ there exists a constant $\delta\in(0,\tfrac12]$ such that
\[
\Pr\bigl(\operatorname{sign}(g_{k,i})=\operatorname{sign}(\nabla_i f(x_k))\mid x_k\bigr)\ge \tfrac12+\delta .
\]
Equivalently,
\[
\Pr\bigl(\operatorname{sign}(g_{k,i})\neq\operatorname{sign}(\nabla_i f(x_k))\mid x_k\bigr)\le \tfrac12-\delta .
\]
\end{assumption}

\begin{assumption}[Bounded second moment]\label{ass:bounded}
There exists $G>0$ such that
\[
\E\bigl[\|g_k\|_2^2\mid x_k\bigr]\le G^2,\qquad\forall k.
\]
\end{assumption}

\section*{Main Convergence Theorem}
\begin{theorem}[Convergence of SignSGD under median bound]\label{thm:main}
Let $\{x_k\}_{k=0}^{T}$ be generated by \eqref{eq:1} with a constant step size $\alpha>0$.  
Assume \ref{ass:smooth}–\ref{ass:bounded} hold and let $\delta$ be the median constant from \ref{ass:median}.  
Define the averaged gradient norm
\[
\Phi_T\;:=\;\frac{1}{T}\sum_{k=0}^{T-1}\E\bigl[\|\nabla f(x_k)\|_1\bigr],
\]
where $\|\cdot\|_1$ denotes the $\ell_1$ norm.  
If the step size satisfies
\[
0<\alpha\le \frac{\delta}{L},
\]
then
\begin{equation}
\Phi_T\;\le\;\frac{f(x_0)-f^{\inf}}{\alpha\delta\,T}\;+\;\frac{\alpha G^2}{2\delta}\;,
\tag{2}
\end{equation}
where $f^{\inf}=\inf_{x}f(x)$. In particular, choosing $\alpha=\sqrt{\frac{2(f(x_0)-f^{\inf})}{G^2 T}}$ yields
\[
\Phi_T = O\!\left(\frac{1}{\sqrt{T}}\right).
\]
\end{theorem}

\section*{Theoretical Properties}
\begin{itemize}
\item \textbf{Stability.} The bound \eqref{eq:2} guarantees that the expected $\ell_1$ gradient norm cannot diverge; the algorithm remains stable for any $\alpha\le\delta/L$.
\item \textbf{Hyper‑parameter sensitivity.} The convergence rate depends linearly on $\delta$, the probability bias of the sign. Larger $\delta$ (more reliable signs) permits larger step sizes and faster descent.
\item \textbf{Comparison to SGD.} Standard SGD under the same smoothness and variance assumptions attains an $O(1/\sqrt{T})$ rate for the $\ell_2$ gradient norm. SignSGD trades the $\ell_2$ norm for an $\ell_1$ norm but enjoys dramatically reduced communication (one bit per coordinate). The rate \eqref{eq:2} matches the $O(1/\sqrt{T})$ behaviour up to constants.
\end{itemize}

\newpage
\section*{Preliminaries (Definitions and Lemmas)}
\begin{lemma}[Smoothness inequality]\label{lem:smooth}
If $f$ is $L$‑smooth then for any $x,y\in\R^{d}$,
\[
f(y)\le f(x)+\langle\nabla f(x),y-x\rangle+\frac{L}{2}\|y-x\|_2^2 .
\]
\end{lemma}
\begin{proof}
Standard consequence of $L$‑smoothness (integrate the gradient Lipschitz bound). \hfill $\square$
\end{proof}

\begin{lemma}[Sign bias inequality]\label{lem:signbias}
For any scalar $a\in\R$ and random sign $s\in\{-1,+1\}$ satisfying
\[
\Pr(s=\operatorname{sign}(a))\ge \tfrac12+\delta,
\]
it holds that
\[
\E[s]\;a\;\ge\; \delta\,|a|.
\]
\end{lemma}
\begin{proof}
Write $\Pr(s=\operatorname{sign}(a))=p\ge\tfrac12+\delta$, $\Pr(s\neq\operatorname{sign}(a))=1-p\le\tfrac12-\delta$.
If $a\ge0$, then $\E[s]=p-(1-p)=2p-1\ge 2\bigl(\tfrac12+\delta\bigr)-1=2\delta$,
so $\E[s]a\ge 2\delta a= \delta|a|$.
The case $a<0$ is symmetric. \hfill $\square$
\end{proof}

\section*{Main Proof (Step‑by‑Step Derivation)}
We prove Theorem~\ref{thm:main}. Throughout, expectations are taken w.r.t. all randomness up to iteration $k$ unless otherwise noted.

\noindent\textbf{Step 1. Smoothness bound for one iteration.}  
Using Lemma~\ref{lem:smooth} with $x=x_k$, $y=x_{k+1}=x_k-\alpha\,\operatorname{sign}(g_k)$,
\begin{align}
f(x_{k+1}) &\le f(x_k)-\alpha\langle\nabla f(x_k),\operatorname{sign}(g_k)\rangle
               +\frac{L\alpha^2}{2}\|\operatorname{sign}(g_k)\|_2^2 . \tag{3}
\end{align}
Since each coordinate of $\operatorname{sign}(g_k)$ belongs to $\{-1,0,1\}$,
\[
\|\operatorname{sign}(g_k)\|_2^2 = \sum_{i=1}^{d}\mathbf{1}_{\{g_{k,i}\neq0\}}\le d .
\tag{4}
\]

\noindent\textbf{Step 2. Take conditional expectation.}  
Condition on $x_k$ and use unbiasedness \ref{ass:unbiased} together with Lemma~\ref{lem:signbias} coordinate‑wise.
For coordinate $i$,
\[
\E\bigl[\operatorname{sign}(g_{k,i})\mid x_k\bigr]\,\nabla_i f(x_k)
\;\ge\; \delta\,|\nabla_i f(x_k)|.
\tag{5}
\]
Summing over $i$ yields
\[
\E\bigl[\langle\nabla f(x_k),\operatorname{sign}(g_k)\rangle\mid x_k\bigr]
\;\ge\;\delta\|\nabla f(x_k)\|_1 .
\tag{6}
\]

\noindent\textbf{Step 3. Bound the variance term.}  
From Assumption \ref{ass:bounded},
\[
\E\bigl[\|g_k\|_2^2\mid x_k\bigr]\le G^2 .
\tag{7}
\]
Because $|\,\operatorname{sign}(g_{k,i})|\le 1$, we also have $\|\operatorname{sign}(g_k)\|_2^2\le d\le G^2$ after possibly redefining $G$ to dominate $\sqrt d$; for clarity we keep the generic bound $G^2$:
\[
\E\bigl[\|\operatorname{sign}(g_k)\|_2^2\mid x_k\bigr]\le G^2 .
\tag{8}
\]

\noindent\textbf{Step 4. Conditional expected descent.}  
Taking expectation of \eqref{eq:3} conditioned on $x_k$ and applying \eqref{eq:6}–\eqref{eq:8},
\begin{align}
\E\bigl[f(x_{k+1})\mid x_k\bigr]
&\le f(x_k)-\alpha\,\delta\|\nabla f(x_k)\|_1
   +\frac{L\alpha^2}{2}G^2 . \tag{9}
\end{align}

\noindent\textbf{Step 5. Unconditional expectation and telescoping sum.}  
Take full expectation and sum \eqref{eq:9} over $k=0,\dots,T-1$:
\begin{align}
\E[f(x_T)] &\le f(x_0)-\alpha\delta\sum_{k=0}^{T-1}\E\bigl[\|\nabla f(x_k)\|_1\bigr]
            +\frac{L\alpha^2 G^2}{2}T . \tag{10}
\end{align}
Rearranging gives
\[
\alpha\delta\sum_{k=0}^{T-1}\E\bigl[\|\nabla f(x_k)\|_1\bigr]
\le f(x_0)-\E[f(x_T)]+\frac{L\alpha^2 G^2}{2}T .
\tag{11}
\]

\noindent\textbf{Step 6. Lower bound $f(x_T)$ by the global infimum.}  
Since $f(x_T)\ge f^{\inf}$,
\[
\alpha\delta\sum_{k=0}^{T-1}\E\bigl[\|\nabla f(x_k)\|_1\bigr]
\le f(x_0)-f^{\inf}+\frac{L\alpha^2 G^2}{2}T . \tag{12}
\]

\noindent\textbf{Step 7. Divide by $\alpha\delta T$ to obtain the average.}
\[
\frac{1}{T}\sum_{k=0}^{T-1}\E\bigl[\|\nabla f(x_k)\|_1\bigr]
\le \frac{f(x_0)-f^{\inf}}{\alpha\delta T}
   +\frac{L\alpha G^2}{2\delta}. \tag{13}
\]
Recall that $L\alpha\le \delta$ by the step‑size restriction $\alpha\le\delta/L$, which guarantees the right‑hand side is non‑negative. Setting $\Phi_T$ as defined in the theorem yields exactly inequality \eqref{eq:2}.

\noindent\textbf{Step 8. Optimizing the step size.}  
Choosing $\alpha$ to balance the two terms in \eqref{eq:13}:
\[
\alpha^\star = \sqrt{\frac{2\bigl(f(x_0)-f^{\inf}\bigr)}{G^2 T}} .
\tag{14}
\]
Since $\alpha^\star\le\delta/L$ for all sufficiently large $T$ (or by clipping to $\delta/L$), plugging \eqref{eq:14} into \eqref{eq:13} gives
\[
\Phi_T \le \frac{G}{\sqrt{2}\,\delta}\,\frac{1}{\sqrt{T}}+ \frac{G}{\sqrt{2}\,\delta}\,\frac{1}{\sqrt{T}}
= O\!\left(\frac{1}{\sqrt{T}}\right).
\tag{15}
\]
Thus the claimed $O(1/\sqrt{T})$ convergence in expectation holds. $\square$

\section*{Rate Derivation (With Constants)}
From \eqref{eq:13} we can write the explicit constant‑wise bound:
\[
\Phi_T \le
\underbrace{\frac{f(x_0)-f^{\inf}}{\alpha\delta}}_{\text{initial gap term}}
\frac{1}{T}
\;+\;
\underbrace{\frac{L\alpha G^2}{2\delta}}_{\text{variance term}} .
\]
Setting $\alpha = \min\bigl\{\frac{\delta}{L},\;\sqrt{\frac{2(f(x_0)-f^{\inf})}{G^2 T}}\bigr\}$ yields the tightest possible bound under the admissible step‑size region.

\section*{Special Cases}
\begin{enumerate}[label=\arabic*.]
\item \textbf{Deterministic gradients.} If $g_k=\nabla f(x_k)$, then $\delta=1/2$ and $G=\|\nabla f(x_k)\|_2$. The bound reduces to the classic gradient descent rate $\mathcal{O}(1/T)$ for smooth non‑convex functions when measured in $\ell_1$ norm.
\item \textbf{High‑dimensional regime.} When $d$ is large but each coordinate’s sign is correct with a modest bias $\delta$, the convergence rate still scales as $1/\sqrt{T}$ and does not deteriorate with $d$ because the $\ell_1$ norm appears linearly.
\item \textbf{Mini‑batch variance reduction.} Increasing the mini‑batch size reduces $G^2$ (by the law of large numbers), thereby allowing a larger optimal step size $\alpha^\star$ and improving the constant in the $O(1/\sqrt{T})$ term.
\end{enumerate}

\section*{Discussion}
The proof above shows that **SignSGD** achieves the same asymptotic convergence order as full‑precision SGD under a mild median‑bias condition on the stochastic sign. The key technical device is Lemma~\ref{lem:signbias}, which transforms a probabilistic statement about sign correctness into a deterministic lower bound on the expected inner product between the true gradient and the sign vector. This lower bound compensates for the loss of magnitude information and yields a clean descent inequality (9). The result highlights that communication‑efficient optimization is possible without sacrificing theoretical guarantees, provided that the stochastic signs are not completely random.

\end{document}